\subsection{Scope and funnel logic}

This review targets journal literature from 2020 to 2026 on automated reinforcement learning for predictive maintenance in manufacturing and industrial machines. The search began with a strict scope: explicit AutoRL for manufacturing predictive maintenance in peer reviewed journals. That strict scope was too sparse to support a stand alone industrial survey. The evidence base therefore widened in a controlled way to three adjacent layers: generic AutoRL methods that automate algorithm choice, hyperparameter tuning, reward shaping, or training adaptation [Par22, Mus22, Die24, Bai24, Wan22c]; industrial predictive and prescriptive maintenance papers that use RL or DRL but still rely on manually configured pipelines [Ong22, Ong22b, Fen22, Hua20, Zha23d, Hu24, Jim23, Net24]; and maintenance papers that automate only part of the RL workflow, such as transfer, hierarchy, or meta adaptation [Ong22b, Abb23, Che24g].

The key result is not only what was found, but what was not found. The retrieved journal literature does not reveal a mature, manufacturing-focused stream in which AutoRL is explicitly developed and deployed for predictive maintenance. This scarcity is itself a substantive finding and supports the claim that AutoRL for industrial maintenance remains an open research area rather than an established application niche [Par22, Sir23, Zha24b, Gia25, Oro25].

\subsection{Research gap}

The literature supports a precise gap statement. RL for industrial maintenance is no longer novel. Reviews and tutorials show a sizable body of work on predictive, condition-based, and prescriptive maintenance using RL and DRL [Sir23, Zha24b, Mar23, Gho21]. AutoRL is also no longer novel as a methods field. Generic work already covers automated hyperparameter optimization, reward tuning, population-based adaptation, and end-to-end RL pipeline construction [Par22, Mus22, Die24, Bai24, Wan22c]. What remains missing is the junction of these two streams. Industrial maintenance papers still assume substantial manual design effort in state construction, reward engineering, algorithm selection, and training configuration. They rarely turn those choices into an automated search or adaptation problem [Ong22, Ong22b, Fen22, Zha23d, Hu24, Jim23, Net24].

The research gap is therefore not whether RL can improve predictive maintenance. The gap is whether predictive maintenance can be served by an AutoRL stack that reduces expert retuning, improves transfer across assets and plants, and makes industrial RL deployment more reproducible and robust [Par22, Zha24b, Gia25, Oro25].

\subsection{Funnel outcome}

\begin{table}[htbp]
\centering
\small
\begin{tabular}{p{0.08\textwidth} p{0.34\textwidth} p{0.45\textwidth}}
\hline
Step & Scope & Outcome \\
\hline
1 & Explicit AutoRL for manufacturing predictive maintenance in journals from 2020 to 2026 & No clear journal stream was identified in the retrieved set \\
2 & Papers that automate part of the RL workflow for maintenance & Only partial proxies were found, such as transfer learning, hierarchical specialization, and meta adaptation [Ong22b, Abb23, Che24g] \\
3 & Industrial predictive and prescriptive maintenance via RL & A usable journal comparison set was found across IIoT manufacturing, production lines, steel, machine-robot systems, and quality-aware maintenance [Ong22, Fen22, Hua20, Zha23d, Hu24, Jim23, Net24] \\
4 & Generic AutoRL methods for technical grounding & A strong method base exists, but it is not maintenance specific [Par22, Mus22, Die24, Bai24, Wan22c] \\
\hline
\end{tabular}
\end{table}

\subsection{AutoRL for predictive maintenance remains a missing layer in the manufacturing literature}

The journal literature from 2020 to 2026 shows a clear asymmetry. Reinforcement learning for predictive, condition-based, and prescriptive maintenance is now an active manufacturing research area, with applications in IIoT production systems, multistage lines, machine-robot environments, and quality-aware maintenance [Sir23, Zha24b, Mar23, Gho21]. In parallel, Automated Reinforcement Learning has matured as a methods field concerned with automating algorithm selection, hyperparameter optimization, reward shaping, and adaptive training [Par22, Mus22, Die24, Bai24, Wan22c]. What remains largely absent is the overlap between these two streams. Across the retrieved journal evidence, no mature manufacturing-focused line of work was found in which AutoRL is explicitly formulated and deployed for predictive maintenance.

This absence is substantive rather than semantic. The current manufacturing papers show that RL can create industrial value when the problem formulation, reward design, and training configuration are chosen well. In IIoT manufacturing, DRL-based predictive maintenance has been shown to improve manpower and equipment resource decisions [Ong22], while transfer learning with demonstrations reduces training wall time and improves adaptation across related manufacturing settings [Ong22b]. In multistage production systems, RL-based maintenance decisions reduce system cost and improve throughput relative to conventional policies [Fen22]. In serial production lines, DDQN-based preventive maintenance learns nontrivial group and opportunistic maintenance behavior without those rules being manually hard-coded [Hua20]. In hot-rolling machine-robot environments, tailored PPO delivers lower maintenance cost under hybrid fault modes and economic dependency [Hu24]. In scrap-based steel production, DRL-based inspection and maintenance planning outperforms conventional maintenance baselines and supports significant financial savings [Net24].

Yet these same studies also expose the unresolved bottleneck. The dominant workflow remains manually engineered. Researchers still handcraft the state variables, define the reward trade-offs, select the RL family, specify the neural architecture, and tune the training settings by experiment [Ong22, Fen22, Hua20, Zha23d, Hu24, Net24]. Even stronger workflow variants, such as transfer learning with demonstrations [Ong22b], hierarchical interpretable maintenance agents [Abb23], and transformer-driven prescriptive maintenance [Zha23d], automate only parts of the pipeline. They do not yet constitute maintenance-specific AutoRL.

The research gap is therefore sharper than a generic claim that more RL work is needed. The gap is that predictive maintenance lacks an AutoRL layer that can reduce expert retuning, formalize reward and policy search, improve transfer across assets and plants, and make industrial RL deployment more reproducible. In other words, the next problem is not whether RL can help maintenance. The next problem is whether RL for maintenance can be made easier to configure and safer to scale beyond one-off, manually tuned case studies [Par22, Mus22, Die24, Bai24, Wan22c].

\subsection{Manufacturing evidence base used as the nearest comparison set}

\begin{table}[htbp]
\centering
\small
\begin{tabular}{p{0.14\textwidth} p{0.24\textwidth} p{0.18\textwidth} p{0.32\textwidth}}
\hline
Paper & Domain & RL family & Main result for this review \\
\hline
[Ong22] & IIoT manufacturing with machine and manpower allocation & PPO-LSTM & Shows that RL can improve maintenance resource management in a realistic stochastic setting \\
[Ong22b] & IIoT manufacturing with transfer across related environments & Transfer learning with DDQN style value learning & Closest manufacturing proxy to AutoRL because it reduces retraining effort \\
[Fen22] & Multistage production systems & RL with approximate dynamic programming & Shows cost and throughput benefits from RL maintenance decisions \\
[Hua20] & Serial production lines with buffers & DDQN & Demonstrates that DRL can learn group and opportunistic maintenance policies \\
[Zha23d] & Sensor-driven prescriptive maintenance with RUL forecasting & Transformer plus DRL & Integrates prognostics with action recommendation, but not AutoRL \\
[Hu24] & No-wait hot-rolling machine-robot collaboration & PPO & Strong industrial case with clearly reported numeric training settings \\
[Net24] & Scrap-based steel production line & DQN, PPO, TRPO comparison & Strong manufacturing case with algorithm comparison \\
[Jim23] & Vacuum packaging machines with quality degradation & Deep RL & Shows quality-aware maintenance benefits in a manufacturing setting \\
[Abb23] & Interpretable predictive maintenance & Hierarchical DRL & Useful proxy for structured policy design beyond flat RL \\
\hline
\end{tabular}
\end{table}

\subsection{Hyperparameter transparency audit of the core manufacturing papers}

\begin{table}[htbp]
\centering
\small
\begin{tabular}{p{0.12\textwidth} p{0.18\textwidth} p{0.37\textwidth} p{0.23\textwidth}}
\hline
Paper & Reported RL setup & Recoverable numeric settings from paper text & Research implication \\
\hline
[Hu24] & Tailored PPO & 1000 iterations, Adam optimizer, ReLu activation, 1 hidden layer, 16 neurons, policy learning rate 0.001, value learning rate 0.01, GAE smoothing factor 0.95, clip parameter 0.2 & Strong example of a manually tuned industrial RL pipeline with enough detail to replicate \\
[Net24] & DQN with PPO and TRPO comparison & The paper clearly compares DQN, PPO, and TRPO, but only a partial algorithm-level setup was recoverable in the current workspace. The accessible article page confirms the comparison but not a full numeric hyperparameter table & Good transparency at the algorithm-comparison level, but not yet a full reproducible AutoRL-style search record \\
[Ong22b] & Transfer learning with demonstrations plus DDQN style value learning & TLD uses two fully connected hidden layers of 24 neurons each, network updated every 100 episodes. DQN and DDQN baselines use two fully connected hidden layers of size 32 and batch size 64. QL baseline uses learning rate 0.9 and discount factor 0.9. Pretraining beyond 8000 epochs equals 80000 demonstration points and showed diminishing gain & Closest manufacturing bridge toward AutoRL, but automation is limited to transfer and demonstration reuse rather than full pipeline search \\
[Ong22] & PPO-LSTM & PPO-LSTM is explicit, but the accessible paper text used here does not expose a full numeric hyperparameter table & Strong practitioner case, weaker configuration transparency \\
[Fen22] & RL with approximate dynamic programming & The accessible text used here does not expose a numeric hyperparameter table & Useful evidence of value, limited evidence for reproducibility burden \\
[Hua20] & DDQN & DDQN is explicit, but the accessible text used here does not expose the numeric tuning table & Important case for production-line maintenance, but not enough detail for a full audit in this workspace \\
[Zha23d] & DQN, PPO, and SAC comparison in a Transformer-driven framework & The paper reports that DQN, PPO, and SAC are compared and that reward weights balance RUL, cost, and downtime. It also exposes a 10-state, 3-action MDP design, but several numeric hyperparameters are not machine-readable in the accessible text used here & Strong evidence that maintenance performance depends on multiple handcrafted choices, yet the search remains manual \\
[Jim23] & Deep RL for condition-based maintenance & Deep RL policy is explicit, but the accessible text used here does not expose a numeric hyperparameter table & Good application evidence, limited tuning transparency \\
[Abb23] & Hierarchical specialized DRL & Hierarchical design is explicit, but the accessible text used here does not expose a numeric hyperparameter table & Useful for structured policy design, not for full hyperparameter audit \\
\hline
\end{tabular}
\end{table}

\subsection{Manual RL pipeline versus a proposed AutoRL maintenance stack}

Current manufacturing papers implicitly follow a manual RL workflow. Domain experts choose the state variables, specify cost and downtime trade-offs, select the algorithm family, and tune the training loop case by case [Ong22, Hua20, Zha23d, Hu24, Net24]. This workflow is sufficient for proving that RL can work, but it does not scale well across plants, machine classes, or changing operating regimes.

A maintenance-specific AutoRL stack would preserve domain knowledge while shifting the expensive search burden away from repeated manual tuning.

\begin{table}[htbp]
\centering
\small
\begin{tabular}{p{0.22\textwidth} p{0.33\textwidth} p{0.33\textwidth}}
\hline
Pipeline stage & Current manufacturing practice & Proposed AutoRL maintenance stack \\
\hline
Problem formulation & Handcrafted MDP or POMDP per case study [Fen22, Hua20, Hu24] & Template library for common maintenance settings such as single machine, line, fleet, and machine-robot systems \\
State design & Manual selection of degradation, RUL, buffer, manpower, and quality variables [Ong22, Jim23, Hu24] & Automated candidate-state generation and ablation guided by task performance \\
Reward design & Manual weighting of uptime, cost, downtime, and quality [Zha23d, Hu24] & Reward-weight search or adaptation over maintenance-specific objectives [Die24] \\
Algorithm choice & Researcher chooses DQN, DDQN, PPO, SAC, TRPO, or custom variants [Hua20, Zha23d, Net24] & Automated model selection under maintenance constraints using an AutoRL controller [Par22, Mus22] \\
Hyperparameter tuning & Manual trial-and-error, often weakly reported [Ong22, Fen22, Hua20] & Bayesian optimization, population-based training, or meta-gradient adaptation [Par22, Bai24, Wan22c] \\
Transfer across assets & Retrain or partially reuse prior models [Ong22b] & Transfer-aware AutoRL with demonstrations, priors, or meta-learned initialization [Ong22b, Che24g] \\
Reproducibility & Benefits often clearer than full configuration details & Full configuration logging, search traces, and uncertainty-aware reporting \\
\hline
\end{tabular}
\end{table}

The value of this comparison is strategic. The manual pipeline already proves that RL is useful for industrial maintenance. What is missing is the automation layer that would make those successes portable, reproducible, and less dependent on expert retuning. That is the cleanest justification for maintenance-focused AutoRL research.

\subsection{AutoRL method base that has not yet been translated well into maintenance}

\begin{table}[htbp]
\centering
\small
\begin{tabular}{p{0.16\textwidth} p{0.30\textwidth} p{0.36\textwidth}}
\hline
Paper & Main automation target & Why it matters for maintenance \\
\hline
[Par22] & Taxonomy of AutoRL including HPO, meta learning, reward design, and architecture search & Gives the clearest map of what an industrial AutoRL maintenance pipeline would need to automate \\
[Mus22] & End-to-end AutoRL pipelines for online and offline RL & Useful blueprint for maintenance pipelines that must support different data regimes \\
[Die24] & Joint optimization of hyperparameters and reward shape & Highly relevant because maintenance rewards often mix downtime, cost, risk, and quality \\
[Bai24] & Population-based hyperparameter adaptation & Relevant for nonstationary plant conditions and long training cycles \\
[Wan22c] & Bayesian generational population-based training & Useful for joint architecture and hyperparameter search in large maintenance state spaces \\
\hline
\end{tabular}
\end{table}

\subsection{Discussion}

The hyperparameter audit reinforces the main gap statement. Even the strongest manufacturing papers remain centered on manually configured RL workflows. Some studies provide enough numeric detail to show the burden clearly, such as the PPO configuration in [Hu24] and the transfer-learning architecture choices in [Ong22b]. However, across the core set, there is still no manufacturing paper that treats algorithm family selection, reward design, and hyperparameter search as a unified automated problem. This supports the claim that AutoRL for predictive maintenance is not yet a mature application stream, but an open systems problem at the junction of industrial maintenance and automated decision learning.

\subsection{Key takeaways}

\begin{itemize}
\item Manufacturing predictive maintenance already has a credible RL evidence base [Ong22, Fen22, Hua20, Hu24, Net24].
\item Explicit AutoRL for manufacturing predictive maintenance is still largely absent [Par22, Sir23, Zha24b].
\item Existing manufacturing papers show value from RL, but still rely on manual state design, reward engineering, algorithm choice, and tuning [Ong22, Hua20, Zha23d, Hu24].
\item Partial bridges exist through transfer learning, hierarchy, and interpretable specialization [Ong22b, Abb23, Che24g].
\item The most defensible research gap is therefore the missing automation layer that would make maintenance RL easier to configure, transfer, and reproduce across industrial settings [Par22, Mus22, Die24, Bai24, Wan22c].
\end{itemize}